\section{模型的建立与求解}

\subsection{问题一模型的建立与求解}

\subsubsection{事件—项目两层数据结构}

设原始明细记录为$k$。先统一患者、开单、报告、项目、医生和设备字段，时间精确到秒。对同一来源、同一患者编号和同一开单时间的记录聚合为事件$i$：
\begin{equation}
i=(\mathrm{source},\mathrm{patientID},\mathrm{orderTime}).
\end{equation}
若医嘱字段形如$[a,\text{彩超}]\,[b,\text{彩超}]$，则按方括号边界拆成两个项目任务，而不按普通逗号拆分，避免把“双侧下肢动脉（股总、股浅、股深等）”错误拆成多个项目。拆分后删除“床旁B超检查加收”“超声计算机图文报告”等非扫查项；床旁字样仍保存在事件和项目标记中。

一个事件的全部有效项目构成集合$\J_i$。其名义工作量和复杂病例工作量分别为
\begin{align}
W_i^{(0)}&=\sum_{j\in\J_i}\tau_{ij}^{(0)},\\
W_i^{(1)}&=\sum_{j\in\J_i}\tau_{ij}^{(1)}.
\end{align}
其中$\tau^{(0)}$取题给区间中点，$\tau^{(1)}$取上界。连续模拟包含167\,361个事件、328\,950个项目任务，具体范围见表\ref{tab:data-scope}。住院在留出边界前保留258个预热事件，门诊、体检留出期背景事件为96\,046个。

\begin{table}[H]
\centering
\caption{连续模拟与留出期数据规模}
\label{tab:data-scope}
\begin{tabular}{lrrr}
\toprule
患者来源 & 连续模拟事件数 & 项目任务数 & 留出期评价事件数 \\
\midrule
住院 & 71,029 & 141,429 & 70,771 \\
门诊 & 55,512 & 69,756 & 55,226 \\
体检 & 40,820 & 117,701 & 40,820 \\
\midrule
合计 & 167,361 & 328,886 & 166,817 \\
\bottomrule
\end{tabular}

\end{table}

\subsubsection{科室需求与多项目结构}

临床科室需求量按训练期申请事件统计。对同一事件包含多个科室字段的极少数情况，以源记录顺序中的首个有效检查项目归属科室；未知科室单独保留，不向高频科室摊派。表\ref{tab:p1-structure}中的科室集中度取自训练期，多项目比例取自全观察期。住院需求分散于87个已知科室，前6科室只占35.45\%；门诊前6科室占62.95\%，适合按重点科室的星期节律预留普通检查容量；体检数据只对应一个体检科室，但其批次规模变化大，不能用“单一科室”等同于“需求稳定”。

\begin{table}[H]
\centering
\caption{训练期科室需求集中度与全观察期多项目事件比例}
\label{tab:p1-structure}
\begin{tabular}{lrrrr}
\toprule
来源 & 训练期事件数 & 已知科室数 & 前6科室占比 & 多项目事件占比 \\
\midrule
住院 & 297,656 & 87  & 35.45\% & 51.15\% \\
门诊 & 245,614 & 133 & 62.95\% & 19.67\% \\
体检 & 239,569 & 1   & 100.00\% & 67.31\% \\
\bottomrule
\end{tabular}

\end{table}

住院和体检的多项目事件比例分别为51.17\%和68.99\%，明显高于门诊的19.72\%。因此设备配置不能只按单项目数量配额：体检批次虽可集中安排，却更需要能够覆盖多个常见部位的共同设备；住院患者则需同时考虑专项能力和减少往返。题面没有给出同次检查固定节约比例，主模型仍逐项目累加题给工时。全观察期边界分析中，若将多项目事件总工时统一乘0.9或0.8，住院平均名义工作量会由17.39分钟降至16.11或14.83分钟；这种节约缺乏项目组合层面的测时依据，故只用于说明容量范围，不进入排程参数。

训练期项目名称按互斥优先级分类。产科III/IV级、产科I/II级和心脏项目使用题给专用时长；血管、腹部/胃肠、泌尿/经腹妇科、经阴道/腔内、浅表器官、儿科专项、介入/定位和一般其他采用题给其他项目区间。先判别专项和检查路径，再判别一般部位，可避免“胎儿心脏”被普通心脏规则覆盖、“经阴道妇科”被经腹妇科覆盖。最终参数见表\ref{tab:duration}。

\begin{table}[H]
\centering
\caption{服务项目时长与训练期相邻报告代理}
\label{tab:duration}
\begin{tabular}{lrrrr}
\toprule
项目族 & 下界/min & 名义/min & 上界/min & 训练期代理样本数 \\
\midrule
产科III/IV级 & 30 & 35 & 40 & 13,451 \\
产科I/II级 & 15 & 17.5 & 20 & 20,095 \\
心脏 & 10 & 12.5 & 15 & 141,017 \\
介入/定位 & 5 & 7.5 & 10 & 6,782 \\
血管 & 5 & 7.5 & 10 & 59,006 \\
经阴道/腔内 & 5 & 7.5 & 10 & 21,244 \\
泌尿/经腹妇科 & 5 & 7.5 & 10 & 75,137 \\
腹部/胃肠 & 5 & 7.5 & 10 & 99,719 \\
浅表器官 & 5 & 7.5 & 10 & 67,825 \\
儿科专项 & 5 & 7.5 & 10 & 3,950 \\
一般其他 & 5 & 7.5 & 10 & 3,400 \\
\bottomrule
\end{tabular}

\end{table}

\subsubsection{需求趋势与峰平阈值}

对来源$c$和自然日$d$，定义申请事件量
\begin{equation}
N_{c,d}^{\mathrm{ord}}
=\left|\left\{i:\mathrm{source}_i=c,\ \lfloor p_i\rfloor=d\right\}\right|,
\end{equation}
报告事件量$N_{c,d}^{\mathrm{rep}}$则以事件报告日期的众数归日。报告量可反映历史完成节奏，但问题二、三的未来需求只能由申请量驱动。

2019--2024年住院开单日均值持续上升：2020年为146.12例/日，2021年161.47例/日，2022年170.47例/日，2023年178.02例/日，2024年194.44例/日。体检呈明显批次性和年份波动，门诊相对平稳。图\ref{fig:monthly}显示三类需求在星期与月份结构之外还存在批量体检与偶发需求集中，因而仅用总体均值不足以设计容量。

\begin{figure}[H]
\centering
\includegraphics[width=\textwidth]{fig01_monthly_demand.pdf}
\caption{三类患者月度日均开单事件与留出边界}
\label{fig:monthly}
\end{figure}

所有运行状态阈值只由训练期确定。对训练期完整日需求序列，定义平峰区间为$[Q_{0.45},Q_{0.55}]$，高峰阈值为$Q_{0.90}$，压力阈值为$Q_{0.95}$：
\begin{align}
\mathcal D_{\mathrm{flat}}&=\{d:Q_{0.45}\le N_d\le Q_{0.55}\},\\
\mathcal D_{\mathrm{peak}}&=\{d:N_d\ge Q_{0.90}\},\\
\mathcal D_{\mathrm{stress}}&=\{d:N_d\ge Q_{0.95}\}.
\end{align}
代表日取最接近相应阈值且时间最早的日期，保证规则可复现。三类开单合计的$Q_{0.45}$、平峰中位、$Q_{0.55}$分别为388.65、412和434例/日，高峰阈值671例/日，压力阈值797.15例/日。训练期峰日代表为2019年11月6日，共671例；平峰代表为2022年3月2日，共412例。留出期超过训练高峰阈值的天数占13.15\%，超过压力阈值占7.95\%，说明容量压力不是由少量异常值造成。

\subsubsection{时间顺序需求检验}

为考察需求规律能否外推，建立三类简单模型。7日季节朴素为
\begin{equation}
\widehat N_d=N_{d-7}.
\end{equation}
过去8周同星期中位数为
\begin{equation}
\widehat N_d=\operatorname{median}\{N_{d-7},N_{d-14},\ldots,N_{d-56}\}.
\end{equation}
日历趋势模型在训练期拟合
\begin{equation}
\begin{aligned}
\log(1+N_d)={}&\beta_0+\beta_1 t_d
+\beta_2\sin\frac{2\pi h_d}{365.25}+\beta_3\cos\frac{2\pi h_d}{365.25}\\
&+\sum_{w=1}^{6}\gamma_w I(w_d=w)
+\sum_{m=2}^{12}\eta_m I(m_d=m)+\varepsilon_d.
\end{aligned}
\end{equation}
其中$h_d$为年内日序。模型用正余弦项刻画年周期，并以岭惩罚控制月份和星期虚拟变量的共线性。预测回到原尺度时使用$\exp(\widehat z)-1$，避免把对数尺度预测直接当作人数。

评价指标为
\begin{align}
\mathrm{WAPE}&=\frac{\sum_d|N_d-\widehat N_d|}{\sum_dN_d},\\
\mathrm{sMAPE}&=\frac1n\sum_d\frac{2|N_d-\widehat N_d|}{|N_d|+|\widehat N_d|}.
\end{align}
在严格时间留出期内，住院开单由过去8周同星期中位数取得最低WAPE 16.04\%，门诊开单为14.04\%；体检开单因批量计划变化，最好的日历趋势模型WAPE仍为69.83\%。三类合计由日历趋势模型取得24.92\%。这说明住院、门诊的星期规律较稳定，而体检量需要未来批次计划才能获得高精度预测。既有医院门诊量研究也观察到趋势、周内效应和序列相关性\cite{zhu2015,xiang2009,luo2017}；后续排程使用实际留出期需求做回放评价，将需求预测误差与调度误差分开。

\begin{table}[H]
\centering
\caption{开单日需求预测的留出期WAPE比较}
\label{tab:forecast-comparison}
\begin{tabular}{lrrr}
\toprule
开单日序列 & 7日季节朴素 & 过去8周同星期中位数 & 日历趋势岭回归 \\
\midrule
住院     & 20.65\% & \textbf{16.04\%} & 16.50\% \\
门诊     & 19.45\% & \textbf{14.04\%} & 22.16\% \\
体检     & 106.43\% & 74.94\% & \textbf{69.83\%} \\
三类合计 & 35.06\% & 25.19\% & \textbf{24.92\%} \\
\bottomrule
\end{tabular}

\end{table}

表\ref{tab:forecast-comparison}说明模型选择按来源分别进行，而不是强制三类需求共用一个预测器。住院和门诊的8周同星期中位数能吸收近期水平变化，又不依赖未来数据；体检的三种误差都较大，根源是团检计划具有外生批次信息。实际部署时，住院、门诊可由历史滚动预测，体检必须接入未来团检名单；若名单不可得，应以高分位压力场景而非单点预测配置容量。

\subsubsection{项目时长与医生效率边界}

历史报告时间不能直接给出扫查时长。数据驱动的超声检查时长估计需要独立的服务过程字段\cite{lai2025}。本文只对同一医生同一报告日的相邻单项目事件计算正间隔，并截取$(0,60]$分钟作为局部吞吐代理。有效单项目代理共510\,957个，中位数7分钟；代理同时包含扫查、书写、等候和批量提交，因此只用于检查题给时长的相对排序。例如产科III/IV级代理中位数21分钟，高于一般项目5--8分钟，支持其更高负荷，但主排程仍使用题给30--40分钟而非21分钟。

对医生$a$在日期$d$按题给名义时长汇总的工作量$L_{ad}$和有效活动窗口$H_{ad}$，构造对数生产率
\begin{equation}
z_{ad}=\log\frac{L_{ad}}{H_{ad}}.
\end{equation}
其中，$H_{ad}$取当日首末报告间隔再加首个事件的名义工时；仅保留不少于3个事件且首末跨度不超过16小时的医生日。对$z_{ad}$作1\%和99\%温莎化后，记医生$a$的样本均值、均值方差和医生间方差分别为$\bar z_a$、$\widehat v_a$和$\sigma_b^2$，采用经验贝叶斯收缩
\begin{align}
\widetilde z_a&=\lambda_a\bar z_a+(1-\lambda_a)\mu_0,\\
\lambda_a&=\frac{\sigma_b^2}{\sigma_b^2+\widehat v_a}.
\end{align}
该指标只描述历史报告吞吐差异，不映射到未来33名医生。主排程中的供给参数来自表6工作日记录：将训练日首末活动时刻重构为活动区间，按星期和5分钟槽统计同时活动人数，主情景取中位数，偏紧情景取下四分位。设经验分位容量为$\widehat C_{w,t}$，题面医生数为33、设备数为22，则实际共享并发容量取
\begin{equation}
C_{w,t}=\min\{\widehat C_{w,t},33,22\}.
\end{equation}
这一双上限反映每项检查同时占用一名医生和一台设备；训练中位容量实际为5--20，如图\ref{fig:doctor-capacity}。

\begin{figure}[H]
\centering
\includegraphics[width=\textwidth]{fig02_doctor_capacity.pdf}
\caption{训练期星期—5分钟共享医生容量}
\label{fig:doctor-capacity}
\end{figure}

\subsubsection{项目—设备兼容模型}

设表1中设备$r$的项目描述集合为$\mathcal E_r$。对规范化项目文本$v_j$，定义匹配类型
\begin{equation}
m(j,r)\in\{0,1,2,3\}.
\end{equation}
其中1表示项目文本直接命中；2表示保留解剖部位和检查路径的同义或包含关系，复合项目取组成部位的同设备交集；3只对应表1明确的广义描述，例如“带有床旁的超声项目”对常规腹部、泌尿和胸水床旁任务的有限覆盖；0表示表1不能确认。主模型兼容集为
\begin{equation}
\R_{ij}=\{r:m(j,r)\in\{1,2,3\}\}.
\end{equation}

对规范化后的211个项目名称逐项匹配表1，其中44个直接命中，60个通过明确的同义或包含关系匹配，92个由题面广义描述覆盖，另有15个无法确认。最后一类任务仍保留在事件分母中，但候选设备集为空。匹配覆盖见图\ref{fig:capability}和表\ref{tab:capability}。

\begin{figure}[H]
\centering
\includegraphics[width=.88\textwidth]{fig03_capability_coverage.pdf}
\caption{项目—设备匹配方式覆盖}
\label{fig:capability}
\end{figure}

\begin{table}[H]
\centering
\caption{项目—设备匹配方式统计}
\label{tab:capability}
\begin{tabularx}{\textwidth}{lrrX}
\toprule
匹配方式 & 项目数 & 项目任务数 & 使用规则 \\
\midrule
直接匹配 & 44 & 185,497 & 表1项目文本直接命中 \\
语义匹配 & 60 & 102,651 & 同义、包含或复合部位交集 \\
明确广义 & 92 & 39,893 & 表1中的广义项目描述覆盖 \\
未确认 & 15 & 909 & 表1无法确认，主模型排除 \\
\bottomrule
\end{tabularx}

\end{table}

床旁条件单独处理。令$b_{ij}=1$表示项目层床旁任务，7号机集合为$\R^{\mathrm{bed}}=\{7\}$，则
\begin{equation}
\label{eq:bedside-capability}
\R_{ij}^{\mathrm{final}}=
\begin{cases}
\R_{ij}\cap \R^{\mathrm{bed}},&b_{ij}=1,\\
\R_{ij},&b_{ij}=0.
\end{cases}
\end{equation}
该定义使“床旁”只改变地点，不改变专项能力。床旁心脏、床旁产科和床旁血管等任务仅在7号机具备相应项目能力时进入候选集合。

\subsubsection{峰平状态下的机器与预约配置}

项目—设备集合给出“哪些机器能够做”，峰平配置进一步回答“何时开放、先安排什么”。表\ref{tab:peak-configuration}以训练期代表日给出三类构成。平峰代表日虽总量为412例，但住院占212例，不能被简单理解为低负荷日；高峰代表日的671例主要来自321例体检批次，调度方式应与住院高峰不同。因而调整以来源结构和专项负荷共同决定，不按总人数同比例放大每台机器。

\begin{table}[H]
\centering
\caption{训练期峰平状态及对应资源配置原则}
\label{tab:peak-configuration}
\begin{tabularx}{\textwidth}{lrrrrX}
\toprule
状态 & 合计/日 & 住院 & 门诊 & 体检 & 机器与预约配置 \\
\midrule
平峰代表日 & 412 & 212 & 161 & 39 & 专项项目固定进入相容诊室；综合项目按最早空闲分配，多项目优先共同设备。 \\
高峰代表日 & 671 & 163 & 187 & 321 & 同时开放数不超过时隙医生容量（最高20）；保留专项机容量，综合机承担可替代项目，体检批次分时进入。 \\
压力状态 & $\ge 797.15$ & -- & -- & -- & 标准班次内先满足政策服务下限；其余背景预约移日或扩展工作时段，不改变设备项目资质。 \\
\bottomrule
\end{tabularx}

\end{table}

设$K_{w,t}^{\mathrm{need}}$为预测任务在星期$w$、时隙$t$所需的并行诊室数，则建议开放数为
\begin{equation}
K_{w,t}^{\mathrm{open}}=\min\{C_{w,t},K_{w,t}^{\mathrm{need}}\}.
\end{equation}
主情景的医生并发上限为20，因此开放设备数随时隙容量变化，不采用22台设备全天同时运行的固定设置。平峰状态下，先开放当日实际有专项需求的心脏、产科、血管、腹部和床旁设备，再由综合设备补足$K_{w,t}^{\mathrm{open}}$；门诊与体检按时段分批，多项目患者优先在共同设备连续完成。高峰状态下，同时开放数可提高至医生容量上限20，专项设备保留给不可替代项目，综合设备承接兼容项目，体检批次错开住院释放集中的时段。压力状态下，现有标准班次先执行既定服务下限；超出能力的可预约背景需求移日，或在医院批准后延长工作时段。三种状态均保持项目设备资质不变，由此形成需求预测、专项资源保留、普通资源弹性开放和背景预约分流相结合的配置方案。
